\documentclass[article,12pt]{elsarticle}
\usepackage[spanish]{babel}
\usepackage[utf8]{inputenc}
\usepackage{amssymb}
\usepackage{flushend}
\usepackage{float}
\usepackage{anysize}
\usepackage{amsmath}
\usepackage{amsfonts}
\usepackage{dcolumn}
\usepackage{amsthm}
\usepackage{caption}
\usepackage{float}
\usepackage{subfigure}
\usepackage{graphicx}
\usepackage{lmodern}
\usepackage{booktabs}
\usepackage{color}
\usepackage{longtable}
\usepackage{rotating}
\usepackage{etoolbox}
\usepackage{comment}
\usepackage{pdflscape}
\usepackage{booktabs}


\begin{document}

\begin{frontmatter}
    \title{Laboratorio de Fenómenos Colectivos \\ Calor Latente de Fusión del Hielo} \\
    \author{Andrés López González$^1$} %Aquí va el nombre de los integrantes que participaron en el informe.
    \address{$^1$ Facultad de Ciencias, Universidad Nacional Autónoma de México, M\'exico CDMX.}
\end{frontmatter}

\section{Materiales}
\begin{itemize}
\item Hielo
\item Báscula digital
\item Termómetro de mercurio
\item Calorímetro de aluminio
\item Agua de laboratorio
\end{itemize}

\section{Ajuste experimental}

Se colocó agua de laboratorio en el interior del calorímetro, registramos con el termómetro de mercurio la temperatura de este así como su masa $m_a$. Después, se colocó un hielo de masa $m_h$ en el agua, inmediatamente se cerró el calorímetro con el termómetro sellando el último hueco de intercambio de calor, lo dejamos disolverse completamente y obtuvimos una temperatura de equilibrio $T_f$, registramos la masa total del agua y la restamos con la masa inicial; para obtener el calor latente investigamos el calor específico del agua de laboratorio y el calorímetro de aluminio, así como pesamos el vasito de calorímetro que estuvo interactuando en el cambio de calor. Nuestra expresión se determinó como,
\begin{align*}
    Q_{perdido\;agua} &= Q_{ganado\;hielo} \\
    -(m_ac_a(T_f-T_a)-m_cc_c(T_f-T_a)) &= m_hc_h(T_f-T_h)+m_hL_h \\
    -(m_ac_a+m_cc_c)(T_f-T_a) &= m_hc_h|T_h| + m_hL_f \\
    \frac{(m_ac_a+m_cc_c)(T_a-T_f)-m_hc_h|T_h|}{m_h} &= L_f
\end{align*}

\newpage
\section{Resultados y análisis}

Mis datos fueron los siguientes \\

\begin{table}[h]
\begin{tabular}{lr}
m.a          & 225.68 \\
T.a            & 22     \\
T.h            & 13     \\
T.f             & 14     \\
m.c        & 45.85  \\
m.h         & 245.64 \\
derretido & 19.96  \\
c.c         & 0.217  \\
c.h         & 0.5    \\
c.a          & 1     
\end{tabular}
\end{table}

De aquí se tiene que el calor latente es el siguiente en calorías y joules \\

\begin{table}[h]
\begin{tabular}{|l|r|}
\hline
Calor latente cal & 87.94066132 \\ \hline
Joules            & 367.943727  \\ \hline
\end{tabular}
\end{table}

\end{document}

